\documentclass[performance]{pol-ren}
\usepackage{gregoriotex}
%% Text of stanzas 1–2, shared by both editions, and the local macros.

%% Diplomatic text of K (back flyleaf of BJ MS 1619).
\newcommand\diplOne{Bogv rodzicza dzewicza bogem slawena maria\\
  Vtwego syna gospodzina matko swolena maria\\
  Sziszczi nam spwczi nam kyrieleyson}
\newcommand\diplTwo{Twego dzela krzcziczela boszicze\\
  Vslisz glosi naplen misli czlowecze\\
  Slisz modlitwǫ yǫsz nosimi\\
  A dacz raczi gegosz prosimi\\
  a naswecze zbozni pobith\\
  posziwocze raski pzbith ky'eleyson}

%% Critical text, normalised after Woronczak (1962).
\newcommand\stanzaOne{Bogurodzica dziewica, Bogiem sławiena Maryja,\\
  u twego syna Gospodzina matko zwolena, Maryja!\\
  Zyszczy nam, spuści nam. Kirielejson.}
\newcommand\stanzaTwo{Twego dziela Krzciciela, Bożyce,\\
  usłysz głosy, napełń myśli człowiecze.\\
  Słysz modlitwę, jąż nosimy,\\
  a dać raczy, jegoż prosimy:\\
  a na świecie zbożny pobyt,\\
  po żywocie rajski przebyt. Kirielejson.}

%% Translation, line by line with the critical text.
\newcommand\transOne{Mother of God, Virgin glorified by God, Mary!\\
  Chosen Mother, Mary, from thy Son the Lord\\
  win for us, send down to us. Lord, have mercy.}
\newcommand\transTwo{For thy Baptist's sake, Son of God,\\
  hear our voices, fulfil the thoughts of men.\\
  Hear the prayer we bring,\\
  and deign to grant what we ask:\\
  on earth a godly life,\\
  after life a dwelling in paradise. Lord, have mercy.}

%% Performance edition: the text as sung, line by line with its translation.
\newcommand\sungOne{Bogurodzica dziewica, Bogiem sławiena Maryja,\\
  u twego syna Gospodzina matko zwolena, Maryja!\\
  Zyszczy nam, spuści nam.\\
  Kyrie eleison.}
\newcommand\sungTwo{Twego dziela Krzciciela, Bożyce,\\
  usłysz głosy, napełń myśli człowiecze.\\
  Słysz modlitwę, jąż nosimy,\\
  a dać raczy, jegoż prosimy:\\
  a na świecie zbożny pobyt,\\
  po żywocie rajski przebyt.\\
  Kyrie eleison.}
\newcommand\sungTransOne{Mother of God, Virgin glorified by God, Mary!\\
  Chosen Mother, Mary, from thy Son the Lord\\
  win for us, send down to us.\\
  Lord, have mercy.}
\newcommand\sungTransTwo{For thy Baptist's sake, Son of God,\\
  hear our voices, fulfil the thoughts of men.\\
  Hear the prayer we bring,\\
  and deign to grant what we ask:\\
  on earth a godly life,\\
  after life a dwelling in paradise.\\
  Lord, have mercy.}




\composer[13th–14th century]{Anonymous}
\title{Bogurodzica}
\edition{First edition}{2026}

\begin{document}
\makemasthead
\gregorioscore[a]{bogurodzica-performance}
\clearpage   % score alone on page 1

\section*{Text and translation}

\begin{textcols}
\stanzanumber{1} & \lines{\sungOne} & \lines{\sungTransOne}\\[0.9em]
\stanzanumber{2} & \lines{\sungTwo} & \lines{\sungTransTwo}\\
\end{textcols}

\begin{rehearsal}
\pnote{The words} Sing the normalised text with modern Polish pronunciation,
the usual practice for this song; the series pronunciation guide gives the
sounds. Sing \pl{Kyrie eleison} in seven syllables, as underlaid, with
\textit{y} as \textit{i}.

\pnote{Pitch} The written final is D. The range is a ninth, C to the D
above, so the score suits most voices as written; transpose freely.

\pnote{Rhythm} The manuscript gives pitches only. Sing in the rhythm of the
words, with notes of roughly equal length. Sing the notes of a group on one
syllable lightly and evenly, and lengthen the last note before each bar
line. The pace is unhurried.

\pnote{Breathing} Breathe at the full bars. The half and quarter bars mark
the rhymes inside each line: a short lift, or a catch breath if needed. The
double bars end each stanza.

\pnote{The B on \pl{spuści}} This B is natural. It is the only B in the song, at
the top of the melody. Some 16th-century copies flatten B elsewhere; this
version has none.

\pnote{The leap at \pl{U twego}} The line jumps an octave from the low D of
\pl{Maryja} to the high D. Later sources smooth it out, but it is in the
oldest copy. Take it cleanly.

\pnote{Forces} The song was sung in church and, by Długosz's account, by
Polish knights before battle, Grunwald 1410 among them. A single voice, a
small schola or a full choir in unison all suit it. One option is a cantor
for each stanza with everyone joining at \pl{Kyrie eleison}.

\pnote{Both stanzas} Stanza 1 is addressed to the Virgin, stanza 2 to
Christ. They share most of their music.

\pnote{The score} Melody and text from the earliest notated copy, Kraków,
Biblioteka Jagiellońska, MS 1619, written shortly after 1400. Text in
normalised Old Polish after J. Woronczak (1962). Bar lines are editorial.
For sources and variants see the critical edition.
\end{rehearsal}
\footcolophon
\end{document}
